% 第 15 章 Error Management Utility
\bichapter{错误管理实用程序}{Error Management Utility}
\label{chap:errmgmt}

本章介绍 IC Validator 工具附带的错误管理实用程序。

错误管理实用程序分以下各节说明：
\begin{itemize}
  \item 错误管理实用程序
  \item 拆分错误数据库
  \item 合并错误数据库
  \item VUE 文件支持
  \item 就地修改数据库
  \item 层次考虑
  \item XML 命令文件
  \item 将错误导出到版图
  \item 管理所引用的 cPYDB 路径
\end{itemize}

% ============================================================
\section{错误管理实用程序}
\subsection*{Error Management Utility}

使用错误管理实用程序（\cmd{icv\_pydb}）可以：
\begin{itemize}
  \item 将单次 IC Validator 运行的错误数据库（PYDB）拆分为多个数据库以便调试。
  \item 从输入 PYDB 选择错误，并仅用所选错误创建新 PYDB。
  \item 将多个 PYDB 数据库合并在一起。
  \item 通过重命名违例或单元、或对错误分类，来更改现有 PYDB 数据库。
  \item 将 PYDB 数据库中的错误导出到版图。
  \item 对 DRC 错误分组。
\end{itemize}

\cmd{icv\_pydb} 实用程序提供命令行选项，可按单元（单独或层次）、违例注释、违例名、错误分类或单元级包围盒进行任意选择。

该实用程序位于 IC Validator 安装目录的 \cmd{bin} 下。表~\ref{tab:icv-pydb-opts} 定义命令行选项。

\begin{longtable}{@{}>{\raggedright\arraybackslash\ttfamily}p{0.36\textwidth} p{0.56\textwidth}@{}}
\caption{错误管理实用程序命令行选项}\label{tab:icv-pydb-opts}\\
\toprule
\textbf{\textrm{命令行选项}} & \textbf{\textrm{定义}} \\
\midrule
\endfirsthead
\caption[]{错误管理实用程序命令行选项（续）}\\
\toprule
\textbf{\textrm{命令行选项}} & \textbf{\textrm{定义}} \\
\midrule
\endhead
\bottomrule
\endfoot

\multicolumn{2}{l}{\textbf{\textrm{输入选项}}} \\
\midrule
-i \textit{input db} &
  指定输入数据库路径，例如 \cmd{run\_details/pydb/PYDB\_top}。必需。此选项必须出现在适用于它的任何输入选择选项之前。所有其他输入选项仅适用于命令行上位于其前的那个输入数据库。 \\
-svc \textit{violation\_comment} / -select\_violation\_comment \textit{violation\_comment} &
  仅包含匹配给定模式的违例注释。 \\
-uvc \textit{violation\_comment} / -unselect\_violation\_comment \textit{violation\_comment} &
  不包含匹配给定模式的违例注释。 \\
-select\_cell \textit{cell} &
  仅包含匹配给定模式的单元。此命令行选项不包含匹配单元下方任意层次上放置的任何单元。 \\
-unselect\_cell \textit{cell} &
  不包含匹配给定模式的单元。此命令行选项不排除匹配单元下方任意层次上放置的任何单元。 \\
-select\_cell\_hierarchical \textit{cell} &
  仅包含匹配给定模式的单元。此命令行选项也包含匹配单元的全部后代。 \\
-unselect\_cell\_hierarchical \textit{cell} &
  不包含匹配给定模式的单元。此命令行选项也排除匹配单元的全部后代。 \\
-select\_interacting \textit{cell,x1,y1,x2,y2} &
  仅包含给定单元中与给定窗口相交的错误。此选项为单元级，不适用于给定单元层次下方的任何错误。窗口选择会合并，使得在某单元中被任一窗口选择包含、且不匹配任何取消选择的错误被包含。 \\
-unselect\_interacting \textit{cell,x1,y1,x2,y2} &
  排除给定单元中与给定窗口相交的错误。此选项为单元级，不适用于给定单元层次下方的任何错误。 \\
-select\_enclosed \textit{cell,x1,y1,x2,y2} &
  仅包含给定单元中被给定窗口包围的错误。此选项为单元级，不适用于给定单元层次下方的任何错误。窗口选择会合并，使得在某单元中被任一窗口选择包含、且不匹配任何取消选择的错误被包含。 \\
-unselect\_enclosed \textit{cell,x1,y1,x2,y2} &
  排除给定单元中被给定窗口包围的错误。此选项为单元级，不适用于给定单元层次下方的任何错误。 \\
-select\_classified Waive,Ignore,Error,... &
  仅包含具有给定分类的错误。 \\
-unselect\_classified Waive,Ignore,Error,... &
  不包含具有给定分类的错误。 \\
-sln \textit{layer\_name} / -select\_layer\_names \textit{layer\_name} &
  仅指定源自匹配给定模式的 runset 层名的错误。仅当 IC Validator 以 \cmd{-vue} 命令行选项运行时，层信息才可用。 \\
-uln \textit{layer\_name} / -unselect\_layer\_names \textit{layer\_name} &
  指定不包含源自给定层的错误。仅当 IC Validator 以 \cmd{-vue} 运行时，层信息才可用。 \\
-sl / -select\_layers \textit{layernum[-range][,datatype[-range]] ...} &
  仅指定源自给定层的错误。对 OpenAccess 版图，层由版图层与 purpose 名指定，而非层与 datatype 号。仅当以 \cmd{-vue} 运行时层信息才可用。 \\
-ul / -unselect\_layers \textit{layernum[-range][,datatype[-range]] ...} &
  指定不包含源自给定层的错误。对 OpenAccess 版图，层由版图层与 purpose 名指定。仅当以 \cmd{-vue} 运行时层信息才可用。 \\
-set / -select\_explorer\_tiers \textit{tier\_number[-range] ...} &
  仅包含属于给定 Explorer tier 编号的错误。 \\
-uet / -unselect\_explorer\_tiers \textit{tier\_number[-range] ...} &
  不包含属于给定 Explorer tier 编号的错误。 \\
-svn / -select\_violation\_name \textit{violation\_name} &
  仅包含匹配给定模式的违例名。违例名对应 runset 中 \cmd{@=} 的左侧。 \\
-uvn / -unselect\_violation\_name \textit{violation\_name} &
  不包含匹配给定模式的违例名。违例名对应 runset 中 \cmd{@=} 的左侧。 \\
-exclude\_heat\_map &
  不将热图复制到输出数据库。若输入与输出数据库中同一单元的热图网格不对齐，可用此选项以避免错误。 \\
-exclude\_error\_classification &
  不将错误分类复制到输出数据库。 \\
-rename\_cells \textit{rename\_cells\_file} &
  根据给定 CSV 文件更改输出中的单元名，第一列为旧单元名，第二列为新单元名。若未指定输出数据库，可用于就地修改输入数据库。就地修改时，要求无其他用户连接到该数据库。 \\
-rename\_violations \textit{rename\_violations\_file} &
  根据给定 CSV 文件更改输出中的违例注释，第一列为旧注释，第二列为新注释。若未指定输出数据库，可用于就地修改输入数据库。就地修改时，要求无其他用户连接到该数据库。 \\
-new\_top\_cell \textit{orig\_cell,new\_cell[,x,y,angle,reflected,[flat\_placement\_count]]} &
  指定与源数据库不同的目标顶层单元，并可选地提供从源顶层单元到目标顶层单元的变换。若新顶层单元在层次中位于原顶层单元之下，则不需要变换。若新单元在层次中位于原顶层单元之上，则需给出原单元到新顶层单元的单实例变换。源数据库到顶层的实例变换会更新到新顶层单元上下文，以便在新顶层单元中高亮违例。可选地，可提供旧顶层单元在新顶层单元内的平坦放置计数，以在新顶层单元上下文中保持平坦错误计数准确。若未提供平坦放置计数，默认为一次放置。 \\
-classify \textit{classify file} &
  根据给定输入文件对输入数据库中的错误分类。此选项仅可用于就地修改 PYDB。更多信息见``分类文件格式''一节。 \\
-classify\_by\_window {[}interacting\textbar{}enclosed{]} &
  指定如何根据分类文件中的窗口坐标选择错误以进行分类。默认为 \cmd{interacting}。
  \cmd{interacting}：与窗口坐标有任何相交的错误被选中分类。
  \cmd{enclosed}：完全被窗口坐标包围的错误被选中分类。所选错误可触及包围窗口。 \\
-report\_user\_comments All \textbar{} most\_recent &
  指定 JSON waiver 报告中包含哪些用户注释。默认为 All。 \\
-set\_waiver\_context \textit{cellname} \textbar{} ALL &
  用 CSV 文件分类错误时，将给定单元设为正在分类错误的 waiver 上下文。这意味着仅当该单元是正在检查设计的顶层单元时才应用 waivers。这通常是仅适用于 block 级的 waivers 在 block 级运行中的顶层单元。若指定 ALL，则始终应用 waivers。 \\
-set\_waiver\_properties \textit{user\_property\_file.json} &
  从给定 JSON waiver 属性文件设置 waiver 用户属性。此操作仅可用于就地修改 cPYDB 或 OASIS waiver 数据库。 \\
-waiver\_property\_definition\_file \textit{property\_definition\_file.json} &
  指定用于验证 waiver 用户属性的 JSON 文件。覆盖环境变量 \cmd{ICV\_WAIVER\_PROPERTY\_DEFINITION\_FILE}。见``自定义 Waiver 属性''。 \\
-select\_waiver\_properties &
  过滤给定 waiver 数据库，仅包含与给定 JSON 属性选择文件中指定的 waiver 用户属性匹配的 waivers。此操作仅可用于就地修改 cPYDB 或 OASIS waiver 数据库。 \\
-rule\_name\_delimiter &
  指定规则名分隔符正则表达式，用于在按规则名而非完整违例注释选择时（例如处理 waiver 属性时）将规则名与违例注释其余部分分开。规则名分隔符默认为 \cmd{/[[:space:]]+:|[[:space:]]*:[[:space:]]/}。 \\
-replace\_user\_ids \textit{new\_user\_id} &
  将输入 cPYDB 或 OASIS waiver 数据库中的全部用户 ID 全局替换为给定值。此操作仅可用于就地修改 cPYDB 或 OASIS waiver 数据库，并要求无其他用户连接到该数据库。 \\
-delete\_classification\_comments &
  删除输入 cPYDB 或 OASIS waiver 数据库中的全部分类注释。此操作仅可用于就地修改，并要求无其他用户连接到该数据库。 \\
-replace\_db\_path \textit{new\_db\_path} &
  将输入 cPYDB 中存储的数据库路径替换为给定值。此操作仅可用于就地修改 cPYDB，并要求无其他用户连接到该数据库。 \\
-rename\_cpydb\_paths \textit{pattern} \textit{new\_path} &
  将任何匹配给定模式的所引用 cPYDB 路径更改为指定的新 cPYDB 路径。仅更新所引用 cPYDB 的路径分量。此选项仅可用于就地修改 PYDB。 \\
-group\_drc\_errors &
  为所选 DRC 错误确定唯一错误组。若违例标记形状与检查值相同（不论方向），则错误放入同一组。仅某些命令有资格分组。此过程替换任何现有分组。更多信息见``DRC 错误分组''。 \\
-minimum\_errors\_per\_group &
  调整构成一组所需的错误阈值。最小阈值为 3。见``DRC 错误分组''。 \\
-group\_same\_orientation &
  将相同方向的错误分在一组。默认不论方向均分组。见``DRC 错误分组''。 \\

\midrule
\multicolumn{2}{l}{\textbf{\textrm{输出选项}}} \\
\midrule
-o \textit{output db} &
  指定输出数据库路径。若未给出此选项，则就地修改输入数据库。若输出数据库不存在则创建；若存在，则除非指定 \cmd{-overwrite}，否则将来自输入数据库的错误合并到其中。所有其他输出命令行选项适用于命令行上位于其前的输出数据库。对此数据库的输出仅来自自上次指定输出数据库以来、在命令行上位于其前的输入数据库。使用 \cmd{-format} 导出到版图时，\cmd{-o} 指定输出版图路径。 \\
-overwrite &
  若输出数据库存在则覆盖。未指定时，默认行为是合并到现有输出数据库，或不存在时创建新输出数据库。 \\
-vue \textit{output VUE file} &
  创建指向输出 PYDB 的给定 VUE 文件。 \\
-vue\_template \textit{input\_vue\_file} &
  使用给定输入 VUE 文件作为模板创建输出 VUE 文件。 \\
-dbu \textit{microns} &
  可选地以微米为单位指定输出数据库的数据库单位值。创建新数据库时使用该值；合并到现有数据库时可用来验证合并数据库的 DBU。未指定时，创建新输出数据库使用输入数据库的 DBU。导出到输出版图时，若给出此参数则用作输出版图的 DBU；若未给出，则使用原始输入版图的 DBU（可能比输入数据库的 DBU 更粗）。 \\

\midrule
\multicolumn{2}{l}{\textbf{\textrm{导出选项}}} \\
\midrule
-format GDSII\textbar{}OASIS &
  表示输出文件指从输入数据库导出的版图，并指定格式。 \\
-layout \textit{original\_layout} &
  指定导出错误到版图时读取设计层次的原始版图。若未给出，版图路径取自 PYDB。无论是否给出此参数，布局导出流程均需要输入版图。 \\
-c \textit{top\_cell} &
  指定版图顶层单元。若未给出，使用输入 PYDB 的顶层单元。 \\
-violation\_map \textit{violation mapping file} &
  指定将规则与命令映射到输出版图中层与 datatype 对的文件。此 CSV 格式文件包含下列列：``Violation Comment'',``Command'',``Layer'',``Datatype''。违例注释与命令列支持通配符。命令列可选，可留空；若存在，则仅当 runset \cmd{file:line:function} 匹配时选择错误。默认情况下，此文件中无条目的错误不输出。若未给出违例映射文件，则为规则分配任意层与 datatype。 \\
-include\_unmapped &
  表示导出错误到版图时包含违例映射文件中无条目的错误。这些错误被分配尚未包含在违例映射文件中的任意层与 datatype。 \\
-violation\_map\_output \textit{output violation mapping file} &
  写出包含导出中所用全部层/datatype 对的违例映射文件。可在后续运行中复用该文件以再现相同的规则层与 datatype 选择。 \\
-command\_property\_number \textit{property number} &
  导出错误到版图时，将命令的 runset \cmd{file:line:function} 信息写到每个输出多边形的给定属性号。这可能显著增大输出版图尺寸。 \\
-density\_property\_number \textit{property number} &
  将密度违例属性以及未匹配密度 waivers 的属性约束写到每个输出多边形的给定属性号。这可能显著增大输出版图尺寸。 \\
-datatype\_per\_classification &
  将不同分类的错误写到不同 datatype。按下表将对应分类状态的值加到输出 datatype：
  Error/Unclassified=0，Ignore=1，Waive=2，Watch=3，Unmatched=4，Fixed=7，WaiverConflict=9。 \\
-host\_init host1:\#cores ... {[}hostN:\#cores{]} &
  为分布式操作分配主机。 \\

\midrule
\multicolumn{2}{l}{\textbf{\textrm{一般选项}}} \\
\midrule
-list\_cells &
  显示输入数据库中的唯一单元列表。 \\
-list\_cell\_details &
  显示输入数据库中的唯一单元列表，含平坦放置计数与实例变换等详情。 \\
-list\_cell\_descendants &
  显示输入数据库中的唯一单元列表，含每个单元全部唯一后代的详细信息。 \\
-list\_layers &
  若 IC Validator 以 \cmd{-vue} 运行，显示数据库中定义的 assign 层列表。示例见表后。 \\
-list\_violations &
  显示输入数据库中的唯一违例注释列表。 \\
-include\_layers &
  若可用，在违例列表中包含层信息。示例见表后。 \\
-list\_cpydbs &
  显示用于匹配输入 PYDB 中错误分类的 cPYDB 数据库列表。 \\
-list\_waiver\_properties &
  显示给定 cPYDB 或 OASIS waiver 数据库中已定义 waiver 属性的违例与单元对列表，并列出 waiver 属性值。 \\
-check\_waiver\_database &
  检查输入 OASIS waiver 数据库的合规性。 \\
-create\_waiver\_report &
  为输入数据库创建 \cmd{WAIVER\_REPORT} 文件。 \\
-include\_waiver\_details &
  在 \cmd{WAIVER\_REPORT} 文件中包含已 waive 错误的详情。覆盖 runset 设置 \cmd{waiver\_options(report\_waiver\_details)}。 \\
-include\_unmatched\_waivers &
  在 \cmd{WAIVER\_REPORT} 文件中包含未匹配 waivers 的详情。覆盖 runset 的 \cmd{report\_waiver\_details} 设置。 \\
-create\_json\_waiver\_file &
  创建 JSON waiver 详情文件。 \\
-password \textit{password} &
  指定对现有错误分类数据库写访问的口令。指定 \cmd{prompt} 可从命令行读取口令。 \\
-file \textit{command file} &
  使用给定 XML 命令文件控制工具操作。 \\
-help &
  显示用法信息。 \\
-V &
  打印实用程序版本。 \\
\end{longtable}

\cmd{-list\_layers} 示例：
\begin{lstlisting}[style=shell]
icv_pydb -i run_details/pydb/PYDB_AD4FUL -list_layers
...
Layers defined in run_details/pydb/PYDB_AD4FUL:
tox (1/0)
poly (5,0)
well (31,0)
psel (14,0)
cont (6,0)
met1 (8,0)
met2 (10,0)
via (19,0)
\end{lstlisting}

\cmd{-list\_violations} 与 \cmd{-include\_layers} 示例：
\begin{lstlisting}[style=shell]
icv_pydb -i run_details/pydb/PYDB_AD4FUL \
  -list_violations -include_layers
...
Violations in database run_details/pydb/PYDB_AD4FUL:
"Met1 Spacing must be >= 3"
  met1 (8,0)
"Met1 to Met2 spacing must be >= 1 and not touching "
  met1 (8,0)
  met2 (10,0)
"Met2 Spacing must be >= 3.0 and >5.0 for longedges >=50"
  met2 (10,0)
\end{lstlisting}

% ============================================================
\section{DRC 错误分组}
\subsection*{DRC Error Grouping}

错误数据库管理实用程序可将具有相同违例标记形状与属性的 DRC 错误放入组中。错误分组可在 VUE 中使用 Show Unique 与 Highlight Unique 菜单选项，仅显示每组的代表性错误。默认不论方向均进行分组。例如：
\begin{lstlisting}[style=shell]
icv_pydb -i run_details/pydb/PYDB_TOPCELL -group_drc_errors
\end{lstlisting}

错误管理实用程序命令行选项见信息表~\ref{tab:icv-pydb-opts}。

% ============================================================
\section{拆分错误数据库}
\subsection*{Splitting Error Databases}

若设计中有多个 block：
\begin{lstlisting}[style=shell]
              TOP
         /     |       \
    Block1   Block2     Block3
      / \      / \        / \
    A   B      A C        B   D
\end{lstlisting}

结果可拆分为三个错误数据库（PYDB），每个子 block 一个。例如：
\begin{lstlisting}[style=shell]
icv_pydb -i run_details/pydb/PYDB_TOP -select_cell_hierarchical "Block1"\
  -o new_pydb/PYDB_Block1 -vue Block1.vue -vue_template TOP.vue

icv_pydb -i run_details/pydb/PYDB_TOP -select_cell_hierarchical "Block2"\
  -o new_pydb/PYDB_Block2 -vue Block2.vue -vue_template TOP.vue

icv_pydb -i run_details/pydb/PYDB_TOP -select_cell_hierarchical "Block3"\
  -o new_pydb/PYDB_Block3 -vue Block3.vue -vue_template TOP.vue
\end{lstlisting}

使用 \cmd{-select\_cell\_hierarchical} 命令行选项可自动拉入所选单元下方放置的任何后代单元。本例中，命令行选项拉入 Block1 下的 A 与 B、Block2 下的 A 与 C、Block3 下的 B 与 D。
可用 \cmd{-select\_cell} 或 \cmd{-unselect\_cell} 在不包含子树的情况下单独包含或排除单元。或者，可用单条 \cmd{icv\_pydb} 命令指定此操作：
\begin{lstlisting}[style=shell]
icv_pydb -i run_details/pydb/PYDB_TOP -select_cell_hierarchical "Block1"\
  -o new_pydb/PYDB_Block1 -vue Block1.vue -vue_template TOP.vue \
  -i run_details/pydb/PYDB_TOP -select_cell_hierarchical "Block2" \
  -o new_pydb/PYDB_Block2 -vue Block2.vue -vue_template TOP.vue \
  -i run_details/pydb/PYDB_TOP -select_cell_hierarchical "Block3" \
  -o new_pydb/PYDB_Block3
\end{lstlisting}

本例中，每个输出数据库的顶层单元仍为 TOP。要使较低层 block 成为顶层单元，可使用 \cmd{-new\_top\_cell} 命令行选项。
\begin{lstlisting}[style=shell]
icv_pydb -i run_details/pydb/PYDB_TOP -select_cell_hierarchical \
  -new_top_cell TOP,Block1 \
  -o new_pydb/PYDB_Block1 -vue Block1.vue -vue_template TOP.vue
\end{lstlisting}

也允许按违例注释与包围盒选择。例如，要在无单元选择的情况下仅拆出 metal1 规则，可如下运行工具：
\begin{lstlisting}[style=shell]
icv_pydb -i run_details/pydb/PYDB_TOP -svc "M1" \
  -o new_pydb/PYDB_TOP_m1 -vue TOP_m1.vue -vue_template TOP.vue
\end{lstlisting}

但请注意，本例中位于顶层的任何热图不会被复制，因为 TOP 在输出数据库中不再存在。

由于 PYDB 不含单元放置数据，包围盒选择仅在单元内有效，不会沿层次向下查看。例如，当 block1 放置范围为 $(100,200)$--$(300,400)$ 时，要包含 block1 上方的顶层数据，可如下运行：
\begin{lstlisting}[style=shell]
icv_pydb -i run_details/pydb/PYDB_TOP \
  -select_cell_hierarchical "Block1" \
  -select_interacting TOP,100.0,200.0,300.0,400.0 \
  -o new_pydb/PYDB_Block1 -vue Block1.vue -vue_template TOP.vue
\end{lstlisting}

注意，当设置 \cmd{error\_options(flatten\_violations)} 时，在整个层次可用的 IC Validator 运行期间，允许将违例展平到指定 block。将 \cmd{flatten\_violations} 参数与包围盒选择结合使用，可对错误进行地理分区。

也可使用若干命令行选项按错误分类进行选择，以包含或排除某些分类的错误，以及从输出中丢弃错误分类数据。

% ============================================================
\section{合并错误数据库}
\subsection*{Merging Error Databases}

\cmd{icv\_pydb} 实用程序可将多个错误数据库（PYDB）合并为一个，例如当子 block 或不同规则集分别运行时。在合并具有相同顶层单元的多个 PYDB 的简单情况下，可对每个输入数据库运行一次 \cmd{icv\_pydb}。例如：
\begin{lstlisting}[style=shell]
icv_pydb -i metal1run/run_details/pydb/PYDB_TOP -o merged_pydb/PYDB_TOP
icv_pydb -i metal2run/run_details/pydb/PYDB_TOP -o merged_pydb/PYDB_TOP
...
\end{lstlisting}

若 \cmd{merged\_pydb/PYDB\_TOP} 不存在则创建，后续运行将数据合并到现有数据库。或者，可在命令行上组合：
\begin{lstlisting}[style=shell]
icv_pydb -i metal1run/run_details/pydb/PYDB_TOP \
         -i metal2run/run_details/pydb/PYDB_TOP \
         -o merged_pydb/PYDB_TOP
\end{lstlisting}

合并具有不同顶层单元的多个数据库时，提供到共同顶层单元的（单实例）变换，以便 VUE 中的 ``highlight in top cell'' 正常工作。可用 \cmd{-new\_top\_cell} 命令行选项完成。例如：
\begin{lstlisting}[style=shell]
icv_pydb -i block1run/run_details/pydb/PYDB_Block1 \
  -new_top_cell Block1,TOP,100,100,0,0 -o merged_pydb/PYDB_TOP
icv_pydb -i block2run/run_details/pydb/PYDB_Block2 \
  -new_top_cell Block2,TOP,200,100,0,0 -o merged_pydb/PYDB_TOP
\end{lstlisting}

被合并的 PYDB 可能处于不同的 DBU 分辨率。合并到现有 PYDB 时，\cmd{icv\_pydb} 工具始终使用目标数据库的 DBU。不过，可在命令行指定 DBU，以便创建新 PYDB 时使用指定 DBU，合并到现有 PYDB 时可验证现有数据库的 DBU。

\cmd{-dbu} 命令行选项允许以微米为单位的浮点值给出 DBU。若合并到现有数据库且其 DBU 与命令行给定值不同，则发出错误并退出工具。

\begin{noteBox}
合并期间不合并 runset 函数。若同一违例中两个相似的 runset 函数被合并，它们将作为独立函数存在于目标数据库中。
\end{noteBox}

% ============================================================
\section{VUE 文件支持}
\subsection*{VUE File Support}

创建新数据库时（无论通过拆分还是合并），可通过指定 \cmd{-vue} 命令行选项输出指向新数据库的 VUE 文件。也可使用对应的 \cmd{-vue\_template} 命令行选项指定输入 VUE 文件作为模板，仅替换数据库名。否则，相关信息取自输入 PYDB。

% ============================================================
\section{就地修改数据库}
\subsection*{Modifying Databases in Place}

支持某些操作就地修改 PYDB 或 cPYDB，而无需目标数据库。这些操作为：
\begin{lstlisting}[style=shell]
-rename_cells rename_cells_file
-rename_violations rename_violations_file
-classify classify.csv
\end{lstlisting}

对就地修改操作，仅用 \cmd{-i} 命令行选项指定输入 PYDB 或 cPYDB，没有输出数据库。

\begin{noteBox}
\cmd{-classify} 命令行选项仅支持 PYDB，不支持 cPYDB。
\end{noteBox}

% ============================================================
\section{层次考虑}
\subsection*{Hierarchy Considerations}

PYDB 为每个唯一单元存储单元后代信息，包括其全部后代的实例变换。这使得拆分数据库时能正确更新到顶层的实例变换，以支持 ``Highlight In Top Cell'' 功能。

常见用法是完全拆出子 block，并使其成为目标 PYDB 的新顶层单元。例如：
\begin{lstlisting}[style=shell]
icv_pydb -i run_details/pydb/PYDB_TOP -o new_pydb/PYDB_Block1 \
  -select_cell_hierarchical "Block1" -vue Block1.vue -vue_template \
   TOP.vue \
  -new_top_cell TOP,Block1
\end{lstlisting}

在此层次示例中，创建以 Block1 为顶层单元的新 PYDB，并选择 Block1 及其下方的全部违例。由于 PYDB 具有全部单元后代信息，当新顶层单元位于原顶层单元之下时不需要变换。若新顶层单元位于原顶层单元之下，且单元选择包含新顶层单元之上的任何内容，则发出错误。

% ============================================================
\section{XML 命令文件}
\subsection*{XML Command File}

XML 命令文件允许使用命令行上可用的任何选择，以及按 ID 选择个别错误。与使用命令行选项类似，XML 文件可为每个源数据库指定多个源数据库与不同选择。文件可 gzip 压缩。

必须遵循全部标准 XML 格式规则。特别要注意：所有参数值必须用引号括起，尖括号、与号或引号等特殊字符必须转义。示例~\ref{ex:xml-cmd-syntax} 展示 XML 命令文件的语法。

\begin{lstlisting}[style=shell,caption={XML 命令文件语法},label={ex:xml-cmd-syntax}]
<?xml version="1.0"?>
<pydb>
  <pydb_operation>
     <input_pydb path="run_details/pydb/PYDB_AD4FUL"
       field_selection_mode="UNION|INTERSECTION">
         <select_cell pattern="cellA*"/>
         <unselect_cell pattern="cellB*"/>
         <select_cell_hierarchical pattern="cellA*"/>
         <unselect_cell_hierarchical pattern="cellB*"/>
         <select_violation_comment pattern="Rule1*"/>
         <unselect_violation_comment pattern="Rule2*"/>
         <select_enclosed cell="cellname" x1="1.0" y1="2.0" x2="3.0"
          y2="4.0"/>
         <unselect_enclosed cell="cellname" x1="1.0" y1="2.0" x2="3.0"
          y2="4.0"/>
         <select_interacting cell="cellname" x1="1.0" y1="2.0" x2="3.0"
          y2="4.0"/>
         <unselect_interacting cell="cellname" x1="1.0" y1="2.0" x2="3.0"
          y2="4.0"/>
         <select_classified>
           <class name="Waive"/>
           ...
         </select_classified>
         <unselect_classified>
           <class name="Waive"/>
           ...
         </unselect_classified>
         <select_layer_names pattern="m1*"/>
         <unselect_layer_names pattern="m2*"/>
         <select_layers range="10-11,0-55"/>
         <unselect_layers range="10-11,0-255"/>
         <select_explorer_tiers range="<range>"/>
         <unselect_explorer_tiers range="<range>"/>
         <rename_cell old="cellA" new="cellB"/>
         <rename_cells file="rename_cells.csv"/>
         <rename_violation old="Rule1" new="Rule2"/>
         <rename_violations file="rename_violations.csv"/>
         <top_cell old="top1" new="top2" dx="1.0" dy="2.0" angle="90"
           reflected="1" flat_placement_count="1"/>
         <!-- Select individual errors by ID -->
         <select_errors command_id="1" cell_id="2">
            <error id="1"/>
            <error id="2"/>
            ...
         </select_errors>
         <!-- No error IDs given == all errors selected -->
         <select_errors command_id="2" cell_id="2"/>
         <!-- No error IDs given, no cell ID given == entire command selected -->
         <select_errors command_id="3"/>
         <!-- Select errors by exact value of a field -->
         <select_errors field="Instance Name" value="I123"/>
         <select_errors field="Distance" value="4.5"/>
         <!-- Select errors by wildcard pattern -->
         <select_errors field="Instance Name" pattern="I4*"/>
         <!-- Select errors based on a range of values -->
         <select_errors field="Distance" min="1.0" max="2.0"/>
         <!-- Selection conditional on another field for commands like density -->
         <select_errors field="Value" min="1.0" max="2.0"
                when_field="Report" is="Area"/>
         <!-- Classify errors -->
         <classify_errors class="Waive" comment="Waive errors in IP block"
         field_selection_mode="UNION|INTERSECTION">
         <!-- Any selection options may go here -->
         </classify_errors>
         <exclude_heat_map/>
         <exclude_error_classification/>
       </input_pydb>
       ...
       <output_pydb path="/path_to_pydb" dbu="0.0001">
         <overwrite/>
         <vue file="output.vue" template="inblock.vue"/>
       </output_pydb>
       <output_layout file="output_layout_file" format="GDSII|OASIS">
         <violation_map file="mapping_file" out="output_file"/>
         <include_unmapped/>
         <input_layout file="input_layout_file" cell="top_cell"/>
         <command_property number="property_number"/>
         <host_init host="host1:#cores"/>
         <host_init host="hostN:#cores"/>
       </output_layout>
     </pydb_operation>
     ...
</pydb>
\end{lstlisting}

XML 命令文件语法定义如下：
\begin{itemize}
  \item \cmd{<pydb>}\\
    命令文件 XML 文档的顶层包围标签。必须有且仅有一个 \cmd{<pydb>} 标签，其余内容均包在其中。
  \item \cmd{<pydb\_operation>}\\
    单次数据库操作（例如拆分、合并或复制）的顶层包围标签。文件中可有多个 \cmd{pydb\_operation} 标签。
  \item \cmd{<input\_pydb path="path/to/pydb" field\_selection\_mode="UNION|INTERSECTION">}\\
    指定输入 PYDB 数据库，与命令行上的 \cmd{-i} 相同。此节内的全部选择与命令行选项仅适用于给定输入数据库。若未指定输出数据库，可就地修改输入 PYDB。单个 \cmd{pydb\_operation} 中可指定多个输入数据库。\\
    可选属性 \cmd{field\_selection\_mode} 控制基于字段值的 \cmd{<select\_errors>} 标签如何组合。默认为 \cmd{UNION}，表示满足任一 \cmd{<select\_errors field...>} 准则的错误均被包含。若指定 \cmd{INTERSECTION}，则仅包含满足全部 \cmd{<select\_errors field...>} 准则的错误。
    \begin{itemize}
      \item \cmd{<select\_layer\_names pattern="m1*"/>}：仅包含匹配给定模式名（允许通配符）的层名。
      \item \cmd{<unselect\_layer\_names pattern="m2*"/>}：不包含匹配给定模式名的层名。
      \item \cmd{<select\_layers range="10-11,0-255"/>}：仅包含匹配给定范围的层。
      \item \cmd{<unselect\_layers range="10-11,0-255"/>}：不包含匹配给定范围的层。
      \item \cmd{<select\_cell pattern="cellA*"/>}：仅包含匹配给定模式的单元。等价于 \cmd{-select\_cell}，可多次指定。
      \item \cmd{<unselect\_cell pattern="cellA*"/>}：不包含匹配给定模式的单元。等价于 \cmd{-unselect\_cell}，可多次指定。
      \item \cmd{<select\_cell\_hierarchical pattern="cellA*"/>}：仅包含匹配给定模式的单元及其后代。等价于 \cmd{-select\_cell\_hierarchical}，可多次指定。
      \item \cmd{<unselect\_cell\_hierarchical pattern="cellA*"/>}：不包含匹配给定模式的单元或其后代。等价于 \cmd{-unselect\_cell\_hierarchical}，可多次指定。
      \item \cmd{<select\_violation\_comment pattern="..."/>}：仅包含匹配给定模式的违例注释。等价于 \cmd{-svc}，可多次指定。
      \item \cmd{<unselect\_violation\_comment pattern="..."/>}：不包含匹配给定模式的违例注释。等价于 \cmd{-uvc}，可多次指定。
      \item \cmd{<select\_enclosed .../>}：在给定单元中，仅包含被给定窗口坐标包围的错误。等价于 \cmd{-select\_enclosed}，可多次指定。
      \item \cmd{<unselect\_enclosed .../>}：在给定单元中，不包含被给定窗口坐标包围的错误。等价于 \cmd{-unselect\_enclosed}，可多次指定。
      \item \cmd{<select\_interacting .../>}：在给定单元中，仅包含与给定窗口坐标相交的错误。等价于 \cmd{-select\_interacting}，可多次指定。
      \item \cmd{<unselect\_interacting .../>}：在给定单元中，不包含与给定窗口坐标相交的错误。等价于 \cmd{-unselect\_interacting}，可多次指定。
      \item \cmd{<select\_classified> ... </select\_classified>}：仅包含匹配给定分类列表的错误。等价于 \cmd{-select\_classified}。
      \item \cmd{<unselect\_classified> ... </unselect\_classified>}：不包含匹配给定分类列表的错误。等价于 \cmd{-unselect\_classified}。
      \item \cmd{<rename\_cell old="cellA" new="cellB"/>}：重命名此输入数据库中的单个单元。可多次指定。
      \item \cmd{<rename\_cells file="rename\_cells.csv"/>}：根据给定 CSV 文件重命名单元。等价于 \cmd{-rename\_cells}。
      \item \cmd{<rename\_violation old="Rule1" new="Rule2"/>}：重命名此输入数据库中的单个违例。可多次指定。
      \item \cmd{<rename\_violations file="rename\_violations.csv"/>}：根据给定 CSV 文件重命名违例。等价于 \cmd{-rename\_violations}。
      \item \cmd{<top\_cell old="top1" new="top2" dx="1.0" dy="2.0" angle="90" reflected="1"/>}：指定与源数据库不同的目标顶层单元，并提供从源顶层单元到目标顶层单元的变换。若合并到具有不同顶层单元的 PYDB，此标签必需。源数据库到顶层的实例变换会更新到新顶层单元上下文，以便在新顶层单元中高亮违例。
      \item \cmd{<select\_errors command\_id="1" cell\_id="2"> ... </select\_errors>}：使用源数据库中的数值 ID 选择要包含的个别错误。若 \cmd{<select\_errors>} 记录内无 \cmd{<error>} 记录，则包含该命令/单元的全部错误。若未指定 \cmd{cell\_id}，则包含该命令的全部错误。此标签可多次指定。
      \item \cmd{<select\_errors field="Instance Name" value="I123"/>}：包含给定字符串、整数或双精度字段等于给定值的错误。字段名应为 \cmd{LAYOUT\_ERRORS} 文件中的列标题。基于字段的选择选项按其所在 \cmd{<input\_pydb>} 或 \cmd{<classify\_errors>} 节的 \cmd{field\_selection\_mode} 属性组合。
      \item \cmd{<select\_errors field="Instance Name" pattern="I4*"/>}：包含给定字符串字段匹配给定通配符模式的错误。
      \item \cmd{<select\_errors field="Distance" min="1.0" max="2.0"/>}：包含给定双精度或整数字段落在给定范围内的错误。
      \item \cmd{<select\_errors field="Value" min="1.0" max="2.0" when\_field="Report" is="Area"/>}：包含满足给定选择准则的错误。添加 \cmd{when\_field} 与 \cmd{is} 属性使此选择以该行中另一值为条件。对 \cmd{density()} 等函数很有用：单个错误可能有许多 Report 与 Value 对，选择必须仅针对特定 Report。上例选择 Area 在 1.0--2.0 范围内的密度错误。
      \item \cmd{<classify\_errors class="Waive" comment="..." field\_selection\_mode="UNION|INTERSECTION"> ... </classify\_errors>}：按给定类别对此输入 PYDB 中的错误分类，并应用可选用户注释。可在 \cmd{<classify\_errors>} 节内指定任何错误选择选项，并与节外选择选项组合。单个 \cmd{<input\_pydb>} 节中可指定多个 \cmd{<classify\_errors>} 节。可选 \cmd{field\_selection\_mode} 控制该节内基于字段的 \cmd{<select\_errors>} 如何组合；默认为 \cmd{UNION}。
      \item \cmd{<exclude\_heat\_map/>}：指定不从输入数据库向输出数据库复制热图数据。等价于 \cmd{-exclude\_heat\_map}。
      \item \cmd{<exclude\_error\_classification/>}：指定不从输入数据库向输出数据库复制错误分类数据。等价于 \cmd{-exclude\_error\_classification}。
    \end{itemize}
  \item \cmd{<output\_pydb path="path/to/pydb" dbu="0.0001">}\\
    指定输出 PYDB 数据库，与命令行上的 \cmd{-o} 相同。\cmd{dbu} 参数可选，等价于 \cmd{-dbu}。单个 \cmd{pydb\_operation} 内最多可指定一次 \cmd{output\_pydb}。
    \begin{itemize}
      \item \cmd{<overwrite/>}：若输出数据库存在则覆盖；否则合并到现有输出数据库。等价于 \cmd{-overwrite}。
      \item \cmd{<vue file="output.vue" template="inlib.vue"/>}：创建指向新输出数据库的给定 VUE 文件。等价于 \cmd{-vue}。可选 \cmd{template} 参数指定用作模板的现有 VUE 文件。
    \end{itemize}
  \item \cmd{<output\_layout file="output\_layout\_file" format="GDSII|OASIS">}\\
    指定将来自输入 PYDB 的违例导出到由给定文件与格式指定的输出版图。
    \begin{itemize}
      \item \cmd{<violation\_map file="mapping\_file" out="output\_file"/>}：为此输出版图指定违例映射文件或输出违例映射文件。
      \item \cmd{<include\_unmapped/>}：若使用违例映射文件，指定映射文件中无条目的违例包含在输出中，位于任意分配的层与 datatype 上。
      \item \cmd{<input\_layout file="input\_layout\_file" cell="top\_cell"/>}：指定从中读取设计层次的输入版图。
      \item \cmd{<command\_property number="property\_number"/>}：指定将 runset \cmd{file:line:function} 以给定属性号添加到输出版图中的每个多边形。这可能显著增大输出版图尺寸。
      \item \cmd{<host\_init host="host1:\#cores"/>}：为分布式操作分配主机。
    \end{itemize}
\end{itemize}

\subsection{用 XML 命令文件分类错误}
\subsubsection*{Classifying Errors With the XML Command File}

XML 命令文件提供使用 \cmd{icv\_pydb} 工具分类错误的方式。可用于其他 \cmd{icv\_pydb} 操作的任何选择选项也可用于选择待分类错误。示例~\ref{ex:common-waiver-xml} 说明若干常见 waiver 场景。

\begin{lstlisting}[style=shell,caption={常见 Waiver 场景},label={ex:common-waiver-xml}]
<?xml version="1.0"?>
<pydb>
  <pydb_operation>
     <input_pydb path="run_details/pydb/PYDB_AD4FUL">

          <!-- Waive all the errors for an entire rule -->
          <classify_errors class="Waive" comment="Waiver for rule M1.R1">
            <select_violation_comment pattern="M1.R1*"/>
          </classify_errors>

          <!-- Ignore all the errors in particular cells -->
          <classify_errors class="Ignore" comment="Ignoring errors in IP cells">
            <select_cell pattern="MYIP*"/>
          </classify_errors>

          <!-- Classify errors in a window -->
          <classify_errors class="Watch" comment="Fix in next rev">
            <select_enclosed cell="TOP" x1="0.0" y1="0.0" x2="10.0"
             y2="5.0"/>
          </classify_errors>

          <!-- Waive metal1 spacing errors based on Distance -->
          <classify_errors class="Waive" comment="Waive these">
            <select_layer_names pattern="metal1*"/>
            <select_errors field="Distance" min="1.0" max="2.0"/>
          </classify_errors>

          <!-- Waive density errors based on Area -->
          <classify_errors class="Waive" comment="Accepted density values">
            <select_violation_comment pattern="DN*"/>
            <select_errors field="Value" min="0.0" max="0.5"
             when_field="Report" is="Area"/>
          </classify_errors>

          <!-- Waive PERC errors based on combination of device name and comment -->
          <classify_errors class="Waive" comment="Waiving M1 path errors"
           field_selection_mode="UNION">
            <select_violation_comment pattern="Rule_1*"/>
            <select_errors field="Device name (Type)" pattern="*/M1 *"/>
            <select_errors field="Comment"
             pattern="Gate has no path to power, ground, or top level port"/>
          </classify_errors>

          <!-- Waive PERC errors with certain instance names -->
          <classify_errors class="Waive" comment="Waiving specific instances">
           <select_violation_comment pattern="Rule_2*"/>
           <!-- The default field_selection_mode is "UNION"-->
           <select_errors field="Reference Instance (Cell Name)"
            pattern="XI123 *"/>
           <select_errors field="Reference Instance (Cell Name)"
            pattern="XI456 *"/>
          </classify_errors>

     </input_pydb>
  </pydb_operation>
</pydb>
\end{lstlisting}

% ============================================================
\section{将错误导出到版图}
\subsection*{Exporting Errors to a Layout}

\cmd{icv\_pydb} 实用程序可将一个或多个错误数据库（PYDB）中的错误导出到 GDSII 或 OASIS 版图文件。支持 \cmd{icv\_pydb} 提供的全部选择选项。必须从用于 IC Validator 运行的原始输入版图读取设计层次。层与 datatype 号可由违例映射文件指定，或任意分配。

违例映射文件为 CSV 文件，含下列列：
\begin{itemize}
  \item \textbf{Violation comment}（必需）：违例或规则注释的完整文本。支持通配符匹配。
  \item \textbf{Command}（可选）：给定规则内特定 IC Validator 命令的 runset \cmd{file:line:function}。支持通配符匹配。若此列为空，则匹配该违例或规则内的全部命令。
  \item \textbf{Layer}（必需）：用于输出匹配违例的层号。
  \item \textbf{Datatype}（必需）：用于输出匹配违例的 datatype 号。
\end{itemize}

若给出违例映射文件，则仅导出映射文件中列出的规则与命令。要包含映射文件中未列出的规则与命令，使用 \cmd{-include\_unmapped} 命令行选项。

若未给出输入版图，\cmd{icv\_pydb} 尝试使用 PYDB 中存储的原始版图路径读取设计层次。

\begin{lstlisting}[style=shell,caption={导出到 OASIS 的示例},label={ex:export-oasis}]
icv_pydb -i run_details/pydb/PYDB_AD4FUL \
      -o errors.oas -format OASIS -layout AD4FUL.GDS -c AD4FUL \
      -violation_map violation_map.csv
\end{lstlisting}

\begin{lstlisting}[style=shell,caption={带选择并创建含所用层与 datatype 的参考违例映射文件},label={ex:export-with-map}]
icv_pydb -i run_details/pydb/PYDB_AD4FUL \
      -svc "Met1*" -select_cell "AD4FUL" \
      -o errors.oas -format OASIS \
      -violation_map_output violation_map_used.csv
\end{lstlisting}

\begin{lstlisting}[style=shell,caption={从 PYDB 创建含任意分配层与 datatype 号的违例映射文件},label={ex:create-viol-map}]
icv_pydb -i run_details/pydb/PYDB_AD4FUL \
      -violation_map_output violation_map.csv
\end{lstlisting}

% ============================================================
\section{管理所引用的 cPYDB 路径}
\subsection*{Managing Referenced cPYDB Paths}

当错误分类数据库（cPYDB）用于在 IC Validator 运行期间应用错误分类时，错误数据库（PYDB）包含指向该 cPYDB 的链接。分类注释与时间戳等部分详情不存储在 PYDB 中，而是在需要时由 VUE 等工具从所引用的 cPYDB 读取。若 cPYDB 不在原位置，这些详情可能缺失。\cmd{icv\_pydb} 实用程序可编辑所引用 cPYDB 的路径，以便 VUE 能加载这些分类详情。

\begin{lstlisting}[style=shell,caption={列出 PYDB 所引用的 cPYDB},label={ex:list-cpydbs}]
icv_pydb -i run_details/pydb/PYDB_AD4FUL -list_cpydbs

...

cPYDBs referenced by run_details/pydb/PYDB_AD4FUL:
  Path: /global/cpydb1
  Name: CPYDB1

   Path: /global/cpydb1
   Name: CPYDB1A

   Path: /global/cpydb2
   Name: CPYDB2
\end{lstlisting}

\begin{lstlisting}[style=shell,caption={使用 -rename\_cpydb\_paths 更新 cPYDB 路径},label={ex:rename-cpydb-paths}]
icv_pydb -i run_details/pydb/PYDB_AD4FUL \
  -rename_cpydb_paths "*cpydb1" "/new/path/cpydb1" \
  -rename_cpydb_paths "*cpydb2" "/new_path/cpydb2"

...

Renamed cPYDBs:
  Old Path: /global/cpydb1
  New Path: /new/path/cpydb1
  Name: CPYDB1

   Old Path: /global/cpydb1
   New Path: /new/path/cpydb1
   Name: CPYDB1A

   Old Path: /global/cpydb2
   New Path: /new/path/cpydb2
   Name: CPYDB2
\end{lstlisting}
